% !TEX program = lualatex
% Auto-generated from YAML (v3) — 04: 情報通信ネットワーク（ネットワークとインターネット）

\documentclass[aspectratio=43,professionalfonts,handout]{beamer}
\usepackage{beamer_template}

\newcommand{\LectureNum}{04}
\newcommand{\LectureTitle}{情報通信ネットワーク（ネットワークとインターネット）}
\newcommand{\CourseName}{情報リテラシー}
\newcommand{\TextPages}{pp.106-109}
\newcommand{\NextTopic}{デジタル通信の仕組みとセキュリティ}
\newcommand{\NextPages}{pp.110-113}

\begin{document}
% --- Slide 1: title ---

\begin{frame}{\CourseName\quad 第\LectureNum 回「\LectureTitle 」}
  {\small \TermName\quad \TeacherName\quad ／\quad 教科書 \TextPages}

  \vfill

  \begin{block}{今日の流れ}
    \begin{enumerate}
      \item 情報通信ネットワークとインターネット
      \item クライアントサーバシステムとWWW
      \item IPアドレスとDNS
      \item インターネットへの接続方法
      \item ルータとNAT
    \end{enumerate}
  \end{block}
\end{frame}

% --- 目標 ---

\begin{frame}{今日の目標}
  \begin{block}{この授業が終わったらできること}
    \begin{itemize}
      \item 情報通信ネットワークの仕組みや構成及びプロトコルの役割や技術を理解し，社会における情報通信ネットワークの役割を説明できる。
    \end{itemize}
  \end{block}
\end{frame}

% --- Slide 2: terms ---

\begin{frame}{基本用語}
  \begin{description}[labelwidth=6em]
    \item[\kw{情報通信ネットワーク}]
      コンピュータやスマートフォンなど，複数の機器を有線や無線の通信回線で接続し，互いに情報通信ができるようにしたもの
    \item[\kw{LAN（Local Area Network）}]
      学校や会社などの組織内にある情報通信ネットワーク。ハブやアクセスポイント，ルータなどのネットワーク接続機器を通して構成される
    \item[\kw{WAN（Wide Area Network）}]
      組織どうしを接続するネットワーク。通常はISP（インターネットサービスプロバイダ）が接続サービスを行っている
    \item[\kw{インターネット}]
      WANどうしが世界規模で相互に接続された情報通信ネットワーク
    \item[\kw{ISP（Internet Service Provider）}]
      インターネットへの接続サービスを提供する事業者。プロバイダともいう
  \end{description}
\end{frame}

% --- Slide 3: terms ---

\begin{frame}{クライアントサーバとWWW}
  \begin{description}[labelwidth=6em]
    \item[\kw{サーバ（server）}]
      インターネットでWebページの閲覧などの際に，それぞれの役割に応じたサービスを提供するコンピュータ
    \item[\kw{クライアント（client）}]
      サーバに対してサービスを要求し，サービスを受ける利用者のコンピュータ
    \item[\kw{クライアントサーバシステム}]
      サーバとクライアントが連携して機能するシステム
    \item[\kw{WWW（World Wide Web）}]
      インターネット上で情報を公開・閲覧するための仕組み。WebページはHTMLで記述され，ハイパーリンクで結びつく
    \item[\kw{HTML（HyperText Markup Language）}]
      Webページを記述するための言語
    \item[\kw{URL（Uniform Resource Locator）}]
      Webページなどの情報が格納されている場所と通信方式を示した文字列
  \end{description}
\end{frame}

% --- Slide 4: points ---

\begin{frame}{URLの構造}
  \begin{block}{ポイント}
    \begin{itemize}
      \item 例：https://www.example.ed.jp/index.html
      \item プロトコル（スキーム）：https（暗号通信のWeb）／http（通常のWeb）／ftp（ファイル転送）
      \item ドメイン名：www.example.ed.jp（サーバ名・組織名・組織種別・国名）
      \item パス名：/index.html（情報が収められているフォルダ名やファイル名）
    \end{itemize}
  \end{block}

  \vfill

  \begin{exampleblock}{補足}
    \begin{itemize}
      \item 組織種別の例：ac（大学・学校法人）/ co（企業）/ go（日本政府機関）/ ed（教育機関・小中高）
    \end{itemize}
  \end{exampleblock}
\end{frame}

% --- Slide 5: terms ---

\begin{frame}{IPアドレスとDNS}
  \begin{description}[labelwidth=6em]
    \item[\kw{IPアドレス}]
      インターネットに接続された機器に割り当てられる固有の番号。IPv4は32ビット，IPv6は128ビット
    \item[\kw{DNS（Domain Name System）}]
      ドメイン名（文字列）とIPアドレス（数値）を対応付ける仕組み。DNSサーバが名前解決を行う
    \item[\kw{ドメイン名}]
      数字で表されるIPアドレスを人間が覚えやすい文字列で表したもの（例：www.kantei.go.jp → 202.214.216.10）
  \end{description}
\end{frame}

% --- Slide 6: points ---

\begin{frame}{インターネットへの接続方法}
  \begin{block}{ポイント}
    \begin{itemize}
      \item ■学校内などのLANに接続：LANケーブルやアクセスポイントを経由して接続する
      \item ■街中でWi-Fi接続：無線LANの環境でアクセスポイントを経由して接続する
      \item ■スマートフォンから接続：携帯電話用の通信回線（基地局）を経由して接続する
      \item ■家庭から接続：ISPと契約して光ファイバなどの接続回線を引き込み，ルータを経由して接続する
    \end{itemize}
  \end{block}

  \vfill

  \begin{exampleblock}{補足}
    \begin{itemize}
      \item フリーWi-Fiにつなげるときはセキュリティに注意しよう
    \end{itemize}
  \end{exampleblock}
\end{frame}

% --- Slide 7: terms ---

\begin{frame}{ルータとNAT}
  \begin{description}[labelwidth=6em]
    \item[\kw{ルータ（router）}]
      ネットワーク間のパケットを中継する機器。複数のネットワークを接続する
    \item[\kw{NAT（Network Address Translation）}]
      プライベートIPアドレスとグローバルIPアドレスを相互変換する仕組み。家庭内の複数端末が1つのグローバルIPアドレスを共有できる
  \end{description}
\end{frame}

% --- Slide 8: exercise ---

\begin{frame}{演習}
  {\small 各自で取り組んでください。}

  \vfill

  \begin{block}{基本問題}
    \begin{itemize}
      \item LANとWANとインターネットの関係を説明せよ
      \item クライアントサーバシステムとは何か説明せよ
      \item DNSの役割を説明し，URLの「www.example.ed.jp」を構成要素に分けて説明せよ
    \end{itemize}
  \end{block}

  \vfill

  \begin{exampleblock}{応用問題（余裕がある人）}
    \begin{itemize}
      \item NATが必要な理由をIPv4アドレスの枯渇問題と関連付けて説明せよ
    \end{itemize}
  \end{exampleblock}
\end{frame}

% --- Slide 9: assignment ---

\begin{frame}{課題・次回予告}
  \begin{importantbox}[課題（次回までに提出）]
    教科書 pp.106-109 を復習すること。次週はデジタル通信の仕組みとセキュリティを学ぶ（pp.110-113）
  \end{importantbox}

  \vfill

  \begin{block}{次回}
    「\NextTopic 」（教科書 \NextPages ）
  \end{block}

  \vfill

  {\small
  質問があれば授業後またはTeamsで受け付けます。}
\end{frame}

% --- チェックリスト ---

\begin{frame}{まとめ・チェックリスト}
  \begin{block}{自己チェック（できたら $\checkmark$ をつけよう）}
    \begin{itemize}
      \item[$\square$] LAN・WAN・インターネットの関係とクライアントサーバシステムを理解した
      \item[$\square$] WWW・URL・IPアドレス・DNSの仕組みを学んだ
      \item[$\square$] インターネット接続方法（学校LAN/Wi-Fi/スマートフォン/家庭）とルータ・NATを把握した
    \end{itemize}
  \end{block}

  \vfill

  {\small 質問があれば授業後またはTeamsで受け付けます。}
\end{frame}

\end{document}
